\section{A Proposed Simulation: Norm Formation under Prestige-Biased Precision}%
\label{sec:simulation}

The formal account in Section~\ref{sec:formal} yields a set of testable predictions about
regime transitions in a community of active-inference agents: as the precision
gain on inter-agent observation channels rises, the community should traverse
states of heterogeneous individual priors, broad consensus, bifurcated
echo-chamber attractors, and finally degenerate fixation on a single agent's
prior. We sketch here a minimum viable simulation that operationalises these
predictions. The base platform is the pymdp epistemic-communities setup of
\citet{albarracin2022EpistemicCommunities}, which already demonstrates that
self-confirming sampling under shared priors produces stable echo chambers
without any malign design. We extend that setup with an explicit
\emph{prestige-cue mechanism} that modulates the precision on inter-agent
observation channels, allowing us to sweep through the predicted regime
sequence in a single parameter.

\subsection*{Setup}

We consider $N \approx 20$ agents, each running a discrete-state active-inference
loop implemented in pymdp. Each agent $i$ maintains a generative model with a
shared discrete outcome space $\mathcal{O}$ and a privately initialised prior
preference vector $\C_i \in \mathbb{R}^{|\mathcal{O}|}$ (the log-preference, in
the sense of \citealt{torresan2025ShapeC}).  Agents observe each other's most
recent actions through a precision-weighted observation channel: the
\emph{effective channel precision} from agent $j$ to agent $i$ is
\begin{equation}
  \prc_{ij} \;=\; \prc_0 \cdot f(\rho_{ij}),
  \label{eq:prestige-precision}
\end{equation}
where $\prc_0 > 0$ is a baseline precision, $\rho_{ij} \in [0,1]$ is a
prestige cue, and $f : [0,1] \to \mathbb{R}_{>0}$ is a monotone gain function
(e.g.\ $f(\rho) = e^{g\rho}$ with gain parameter $g \geq 0$). A higher
$\prc_{ij}$ means that agent $i$ weights agent $j$'s actions more heavily when
updating its own $\C_i$ --- precisely the attentional inheritance mechanism
formalised in Section~\ref{sec:installation}.

\paragraph{Key extension over \citet{albarracin2022EpistemicCommunities}.}
Albarracin et al.\ treat all inter-agent channels with a uniform precision;
the community's attractor structure follows from the topology of which sources
agents prefer to sample. Here we add a \emph{prestige layer} that modulates
channel precision independently of source preference. The prestige cue
$\rho_{ij}$ can be operationalised in two ways that correspond to distinct
cultural-learning regimes identified by \citet{henrichGilWhite2001Prestige}:

\begin{enumerate}[label=(\alph*)]
  \item \textbf{Exogenous prestige.} A small fraction $p_h \approx 0.15$ of
    agents are designated high-prestige at the outset ($\rho_{ij} = 1$ for
    high-prestige $j$, $\rho_{ij} = 0$ otherwise). This recovers the simplest
    case of model-based learning studied by \citet{chudek2012BystanderGaze},
    where bystanders differentially attend to marked individuals.

  \item \textbf{Endogenous prestige.} $\rho_{ij}$ self-organises via a
    Henrich-style success-bias rule: after each round, $\rho_{ij}$ is
    proportional to the reduction in free energy that agent $j$ achieved in the
    preceding episode.  Agents that minimise their prediction error most
    efficiently accrue prestige and attract greater attentional weight,
    implementing the ``copy the successful'' heuristic at the precision level.
\end{enumerate}

Both variants share the functional form~\eqref{eq:prestige-precision}; they
differ only in how $\rho_{ij}$ evolves. Variant (b) is the richer case for
the norm-formation question but variant (a) is analytically cleaner for
isolating the regime transitions.

\subsection*{Figure 4 --- Simulation schematic}

\begin{figure}[ht]
\centering
% NOTE: This TikZ figure uses the `graphs` and `graphdrawing` libraries
% (loaded in main.tex) together with the `force` graph-drawing algorithm.
% These require LuaLaTeX. When compiling with pdfLaTeX, replace the
% \graph{} block with manually positioned \node / \draw commands.
\begin{tikzpicture}[
  agent/.style   = {circle, draw, minimum size=0.75cm, inner sep=1pt,
                    font=\small},
  prestige/.style= {circle, draw, very thick, minimum size=0.75cm, inner sep=1pt,
                    font=\small, fill=gray!20},
  chan/.style    = {->, >=Stealth, gray},
  hichan/.style  = {->, >=Stealth, line width=2.5pt, black!70},
  label distance=2pt
]
% --- manually positioned layout (pdfLaTeX-safe version) ---
% Outer ring: 8 ordinary agents
\foreach \i/\angle in {1/0, 2/45, 3/90, 4/135, 5/180, 6/225, 7/270, 8/315}{
  \node[agent] (a\i) at ({2.4*cos(\angle)},{2.4*sin(\angle)})
        {$\C_{\i}$};
}
% Inner core: 2 high-prestige agents
\node[prestige] (p1) at (-0.7, 0.0) {$\C^{*}_{1}$};
\node[prestige] (p2) at ( 0.7, 0.0) {$\C^{*}_{2}$};

% Low-precision channels (ordinary → ordinary, thin)
\foreach \s/\t in {1/2, 2/3, 3/4, 4/5, 5/6, 6/7, 7/8, 8/1, 1/5, 3/7}{
  \draw[chan] (a\s) -- (a\t);
}
% High-precision channels (ordinary → prestige, thick)
\foreach \s in {1,2,3,4,5,6,7,8}{
  \draw[hichan] (a\s) -- (p1);
  \draw[hichan] (a\s) -- (p2);
}
% Mutual channel between high-prestige agents
\draw[hichan, <->] (p1) -- (p2);

% Prestige labels on thick edges (one representative)
\draw[hichan] (a1) -- node[above right, font=\scriptsize]{$\prc_{ij}\!=\!\prc_0 e^{g\rho_{ij}}$} (p1);

% Legend
\node[agent,  below left=0.3cm and 0.0cm of a8, font=\scriptsize]  (leg1) {};
\node[right=0.05cm of leg1, font=\scriptsize] {ordinary agent ($\C_i$)};
\node[prestige, below=0.45cm of leg1, font=\scriptsize] (leg2) {};
\node[right=0.05cm of leg2, font=\scriptsize] {high-prestige agent ($\C^*$)};
\draw[chan,  below=0.6cm of leg2] (0,-3.6) -- node[right, font=\scriptsize]{low $\prc_{ij}$} ++(0.8,0);
\draw[hichan] (0,-3.95) -- node[right, font=\scriptsize]{high $\prc_{ij}$} ++(0.8,0);
\end{tikzpicture}
\caption{Proposed simulation schematic: $N$ agents with individual prior
preferences $\C_i$ observe each other through precision-weighted channels
modulated by prestige cues $\rho_{ij}$.  High-prestige agents (shaded, inner
ring) attract thick observation channels from all other agents; ordinary
agents are connected by thin channels with baseline precision $\prc_0$.  The
community-level $\C$-distribution evolves via attention-mediated copying as
defined in Eq.~\eqref{eq:prestige-precision}.}
\label{fig:sim-schematic}
\end{figure}

\subsection*{Predicted regime transitions}

As the prestige-precision gain $g$ increases from zero, the model predicts
the following sequence of community-level behaviours. At \emph{low} $g$
($g \approx 0$), inter-agent observation channels carry negligible precision
above the individual's self-model; transmitted information is too weak to
overcome agents' prior heterogeneity, and the population $\C$-distribution
remains broad---no norm forms. At \emph{intermediate} $g$, transmission is
strong enough for iterative updating to draw the community toward a single
high-density region: a shared $\C$ emerges, corresponding to the
\emph{shaped} $\C$ regime of Section~\ref{sec:shape-of-C}. At \emph{high} $g$ the
dynamics bifurcate: sub-communities that share initial $\C$-proximity lock
onto distinct attractors, recovering the echo-chamber result of
\citet{albarracin2022EpistemicCommunities} by a complementary mechanism (there,
echo chambers follow from source-sampling bias; here, they follow from
heterogeneous prestige cue weights). At \emph{very high} $g$, the community
collapses onto a single high-prestige agent's $\C$, producing the
\emph{degenerate fixation} regime: the spiky $\C$ of a community captured by
one model.  These four qualitative outcomes span the spiky/broad/shaped $\C$
taxonomy introduced in Section~\ref{sec:shape-of-C} and predicted by the stability
analysis of Section~\ref{sec:loop}.

\subsection*{Quantities to measure}

The simulation collects three primary observables as functions of $g$:

\begin{enumerate}[label=(\roman*)]
  \item \textbf{Mean pairwise KL divergence} $\bar{D}(g) = \binom{N}{2}^{-1}
    \sum_{i < j} \KL{\C_i}{\C_j}$, measuring aggregate heterogeneity of the
    community's prior-preference distribution; this should be non-monotone in
    $g$, peaking at the echo-chamber bifurcation before collapsing at
    degenerate fixation.

  \item \textbf{Population $\C$-entropy} $H(g) = -\int p(\C) \log p(\C)\,
    \mathrm{d}\C$, estimated by kernel density over the empirical distribution
    of $\C_i$ vectors; a unimodal drop in $H$ marks the consensus regime,
    while a bimodal recovery marks the bifurcation.

  \item \textbf{Attractor count} $K(g)$, estimated by $k$-means clustering on
    the converged $\{\C_i\}$ with silhouette-based model selection; $K$
    should trace the sequence $K > 1 \to K = 1 \to K = 2 \to K = 1$ as $g$
    increases through the four regimes.
\end{enumerate}

These observables admit a complementary spectral-graph reading that
connects the simulation directly to the structural analysis of
\citet{heins2024CollectiveSurprise}. Construct the agent influence graph
$\mathcal{G}$ with edge weights equal to the realised $\prc_{ij}$ at
convergence. The \emph{spectral gap} $\lambda_2(\mathcal{G})$ --- the second
eigenvalue of the normalised Laplacian --- measures how quickly information
diffuses across the weakest bottleneck; a large spectral gap indicates a
well-mixed community (consensus regime), while a small spectral gap flags a
bottleneck consistent with bifurcated echo chambers. The \emph{Fiedler
vector} $\mathbf{v}_2(\mathcal{G})$ identifies the natural partition of the
community: its sign pattern should recover the two attractor clusters at high
$g$. Tracking $\lambda_2$ and $\mathbf{v}_2$ as functions of $g$ thus gives
a spectral read-out of the same regime transitions measured by
$\bar{D}(g)$ and $K(g)$.

\subsection*{Connection to the silly-rules result}

The prediction in regime (ii)---the intermediate-$g$ fragile-consensus
window---has a direct analogue in \citet{koster2022SillyRules}. Silly rules,
in the cooperative-AI setting, are arbitrary behavioural taboos that carry no
direct welfare content but substantially improve agents' capacity to learn
stable enforcement and compliance. In the present framework the mechanism is
transparent: silly rules pump precision on the policy prior
``follow-norms-in-general'' without committing to any specific $\C$-content.
They are, in effect, a device for raising $g$ artificially at the level of
policy selection, making the normative channel more legible without specifying
what travels through it. Our framework predicts that silly rules should be
most beneficial precisely in the intermediate-$g$ regime, where the
consensus-forming process is fragile: the channel precision is high enough to
synchronise agents but the attractor is not yet locked in, making additional
precision-pumping most valuable. At very high $g$ (degenerate fixation) silly
rules add noise around an already-collapsed distribution; at low $g$ they
cannot compensate for an insufficiently precise transmission channel. The
intermediate window is where content-neutral precision amplification does
meaningful normative work, consistent with the \citeauthor{koster2022SillyRules}
finding that spurious normativity enhances acquisition of welfare-relevant
compliance.

